\sectionguard
\section*{The Propers: Interpretive Possibilities}

\begin{studybox}{What this section is}
Everything below is exploratory editorial and AI proposal. None of it is
sourced historical intent, attributed teaching, or documented reception, and
none of it should be quoted as any of those. Each proposal was tested against
the corpus checked for this guide and against a targeted search for its own
distinctive conjunction; the anchors, the search boundary, the novelty result
and the controlling limit of each are recorded in \texttt{research/scope.md}.
Evidence, doctrine, and the literal sense of the appointed texts govern; where
a proposal and a checked witness conflict, the witness wins. Bede's rule binds
here as everywhere in this guide: no allegory may thin out the moral command
the parable ends with.
\end{studybox}

\begin{proposal}{Half-dead is a diagnosis of exactly what the letter can reach}
\pfield{Anchors}{Gospel \cue{Gosp.}, \latin{semivívo relícto} with Augustine's
gloss in \work{Quaestiones Evangeliorum} II.19 --- half alive \latin{ex parte
qua potest intellegere et cognoscere Deum}; Epistle \cue{Ep.}, \latin{líttera
enim occídit}, with Aquinas's mechanism in \work{Super II ad Corinthios}
c.~3, \latin{littera legis dat solam cognitiónem peccáti}.}
\pfield{Mechanism}{Augustine divides the man on the road along a precise line:
he is alive in the part that can know God and dead in the part wasted by sin.
Aquinas, eight centuries later and expounding a different text, describes the
letter's whole competence along the same line: it delivers knowledge of sin
and does not touch the cause. Set the two together and the parable stops being
an illustration of the Epistle and becomes its case history. The letter can
speak to the half of the man that is still awake --- and that half is precisely
what the priest and the Levite, who could read, address by walking past. What
the road needs is an intervention aimed at the other half, and the Samaritan's
first act is not instruction but binding.}
\pfield{Fruit}{``The letter killeth'' stops sounding like a rhetorical
antithesis about covenants and acquires an anthropology: a creature left with
cognition and without power is not helped by more accurate cognition. It also
disciplines the preacher, who is by definition working on the waking half.}
\pfield{What the ordinary element-by-element reading misses}{Read separately,
Augustine's \latin{semivivum} is a remark about original sin and Aquinas's
\latin{occasionaliter} is a remark about Romans 7. Neither witness has the
other's text in view, and the missal's page gives no cue that the two
diagnoses are the same diagnosis.}
\pfield{Strongest limit}{Augustine wrote in about 400 and Aquinas in the
1260s; they are not in conversation, and the shared line is drawn here, not by
them. Augustine's own point at \latin{semivivum} is the survival of the
\latin{imago}, not the failure of written law, and the parable never says the
wounded man had read anything. The proposal is a conjunction of two
anthropologies, not evidence about either author's intention or about the
missal's.}
\end{proposal}

\begin{proposal}{The road supplies two of the three medicines; the rail supplies all three}
\pfield{Anchors}{Gospel \cue{Gosp.}, \latin{infúndens óleum et vinum};
Communion \cue{Comm.}, \latin{et vinum lætíficet cor hóminis: ut exhílaret
fáciem in óleo, et panis cor hóminis confírmet}; and, at one remove, Gradual
\cue{Grad.}, through Cassiodorus's own cross-reference of the Communion's
gladdened face to Ps.~33:6.}
\pfield{Mechanism}{The Samaritan carries oil and wine and nothing else; the
inn is where he stops, pays, and leaves the man in another's keeping. The
Communion antiphon names the same two substances and adds a third the roadside
could not carry --- bread, twice, framing the other two. The tradition has
read each list, separately and in opposite directions: Augustine makes the
Samaritan's oil \latin{consolatio spei bonae} and his wine \latin{exhortatio
ad operandum ferventissimo spiritu}, Gregory makes the wine bite and the oil
soothe, while Cassiodorus reads the Communion's wine as gladdening
\latin{cum sacratum fuerit in sanguinem Domini Christi} and its oil as the
face brightened \latin{cum regale chrisma conficitur}. The proposal is
structural rather than allegorical: the formulary moves the day's medicine
from the road to the table, and the item it gains in transit is the one the
Church can only give indoors.}
\pfield{Fruit}{The Communion stops being a pleasant creation-psalm coda and
becomes the answer to a question the Gospel left open --- what happens at the
inn. It also gives the day's mercy a location: not a gesture on a road, but a
house with provisions and a bill already paid.}
\pfield{What the ordinary element-by-element reading misses}{Element by
element the Gospel's oil and wine are first aid, the Communion's are
providence, and the two lists never meet; the commentary treats each under its
own witnesses because that is where the evidence stops.}
\pfield{Strongest limit}{The orders differ --- Luke has oil then wine, the
antiphon wine then oil --- and the psalm's triad is agricultural, with no
wound anywhere in it. Reading the Samaritan's oil and wine as sacraments is
old and common in devotional commentary on this Sunday, but no checked witness
joins that reading to \textit{this} antiphon, and Cassiodorus's sacramental
reading of the psalm was written with no parable in view. The pairing of
these two chants with this Gospel is attested early in the Gregorian
tradition, which proves association and not purpose.}
\end{proposal}

\begin{proposal}{The day's danger is a verb of thinking, and it survives only in the chant's psalter}
\pfield{Anchors}{Introit \cue{Int.}, \latin{qui cógitant mihi mala}, the chant
reading against the Clementine's \latin{qui volunt mihi mala}; Epistle
\cue{Ep.}, \latin{non quod sufficiéntes simus cogitáre áliquid a nobis};
Gospel \cue{Gosp.}, the lawyer who asks \latin{volens iustificáre seípsum}.}
\pfield{Mechanism}{The Introit's verb is a verified collation finding: the
missal sings \latin{cógitant}, devise, where the Clementine and therefore the
Douay have \latin{volunt}, desire. Augustine's exposition turns on exactly
that reading --- open persecution has passed, ``now there has remained the
malice of them \textit{thinking}.'' The Epistle then disclaims sufficiency at
the same faculty and in the same verb, \latin{cogitáre}; and the Gospel's
antagonist is neither a robber nor a persecutor but a reader engaged in an
interior operation, justifying himself. Across three elements the day's
trouble is located inside the head, and in the Introit it is located there
only because the Mass kept an older psalter than the Bible did.}
\pfield{Fruit}{The formulary's enemies stop being external. A Mass that opens
by asking for rescue from people who \textit{think} evil, argues that we
cannot \textit{think} anything sufficient, and then shows a man ruined by
thinking well of himself, is describing one danger three times --- and the
petition it makes is not for cleverness.}
\pfield{What the ordinary element-by-element reading misses}{The link is
invisible in English: the Douay renders the Clementine, so the Introit reads
``desire evils'' and the verbal chain breaks at its first link. It is also
invisible in any book that prints the Clementine psalm instead of the chant.}
\pfield{Strongest limit}{The Latin words do not mean the same thing.
\latin{Cogitáre} at 2 Cor. 3:5 renders a verb of reckoning or accounting, and
Paul is not describing malice; the Introit's \latin{cógitant} is hostile
intent; and the lawyer's own verbs are \latin{tentans} and \latin{iustificáre},
not \latin{cogitáre} at all. The chant's reading is a fact of Latin psalter
history, not a choice made for this Sunday, and Augustine was expounding
persecutors, not lawyers.}
\end{proposal}

\begin{proposal}{An unanswered cry, then someone else's prayer with the answer printed in it}
\pfield{Anchors}{Alleluia \cue{All.}, Ps.~87:2, \latin{in die clamávi et nocte
coram te}, the opening of the one psalm that never brightens; Offertory
\cue{Off.}, \latin{Precátus est Móyses in conspéctu Dómini}, ending \latin{et
placátus factus est Dóminus}.}
\pfield{Mechanism}{Three chants apart, the Mass sings two prayers with
opposite endings. The first is a cry whose own psalm records no answer at all;
the second is an intercession whose last clause \textit{is} the answer, sung
while the gifts are carried. The difference between them is not intensity but
person: the second prayer is offered by someone else, for people who have
disqualified themselves, and the checked reception says exactly that this is
how such prayers work --- Augustine hears God's \latin{Sine me} as
\latin{dilectio tua in illos intercedit mihi}, your love intercedes with me
for them, so that when our own merits weigh us down we may be relieved by the
merits of those God loves. On that reading the liturgy does not answer the
Alleluia's lament by improving it. It answers it by changing who is speaking.}
\pfield{Fruit}{Intercession stops being a courtesy and becomes the day's
structural solution: the formulary that denies self-supplied sufficiency also
declines to make the sufferer's own petition the instrument of rescue. It also
explains, without sentiment, why a psalm that ends in darkness can be sung
under \latin{Allelúia}.}
\pfield{What the ordinary element-by-element reading misses}{The Alleluia is
usually treated as a bright verse detached from a dark psalm, and the
Offertory as a narrative chant admired for its melody. Their sequence carries
an argument that neither carries alone.}
\pfield{Strongest limit}{Only Ps.~87's opening confession is appointed, so the
psalm's darkness is not in the Mass at all --- the reader supplies it from the
psalter. \latin{Coram te} and \latin{in conspéctu Dómini} are different
phrases, and Theodoret's reading makes the lament exilic Israel's rather than
any individual's. Augustine's remark is on Exodus and has no chant in view,
Alleluia verses are commonly assigned on musical and cyclical grounds, and no
checked witness joins these two chants.}
\end{proposal}

\begin{proposal}{The office named by an etymology is asked for as an effect of Communion}
\pfield{Anchors}{Gospel \cue{Gosp.}, the Samaritan read as \latin{custos},
guardian --- Origen's \latin{sciebat enim se custodem}, Ambrose's
\latin{Samaritani etenim vocabulo custos significatur}, Augustine's
\latin{Samaritanus custos interpretatur}, Gregory's \latin{custos namque
humani generis venerat}; Postcommunion \cue{Postcomm.}, \latin{et páriter
nobis expiatiónem tríbuat, et munímen}.}
\pfield{Mechanism}{Four witnesses across four centuries agree that the
Samaritan's name states an office, and Gregory presses it furthest: charged
with being a Samaritan, the Lord denied the demon and kept silent about the
name, because to deny it would have been to deny the office. The last oration
of the same Mass asks for that office as a sacramental effect. \latin{Munímen}
is what the man on the road lacked and what the inn supplied: he was already
robbed when the guardian arrived, and the last thing the Samaritan does is
guarantee the future --- \latin{quodcúmque supererogáveris, ego \dots\ reddam
tibi}. The Postcommunion's two gifts fall on the same two sides of the story:
\latin{expiátio} for what has already happened on the road, \latin{munímen}
for the road still to come, and \latin{páriter} refuses to rank them.}
\pfield{Fruit}{\latin{Munímen} stops being a vague word for protection and
acquires the parable's content --- a guardian who paid, left instructions, and
promised to return. It also gives Augustine's \latin{adhuc curatur} a
liturgical form: the Mass ends by asking for the continuing treatment the inn
was for.}
\pfield{What the ordinary element-by-element reading misses}{Element by
element the etymology is a curiosity of patristic philology and the
Postcommunion is a compact prayer for grace; neither suggests the other, and
the oration's own transmission history has no Gospel attached to it.}
\pfield{Strongest limit}{\latin{Custos} and \latin{munímen} are unrelated
words with unrelated images --- a person and a fortification --- and the
etymology itself is contested reception rather than settled philology. The
Postcommunion is attested in the Verona collection with no lection beside it,
stands in the Gelasian in a Mass-set the missal's Sundays do not match, and
serves a Lenten weekday in the postconciliar Missal; nothing indicates it was
chosen for this Gospel, and no checked witness joins them.}
\end{proposal}
